\documentclass{beamer}

\def\Tiny{\fontsize{6pt}{6pt}\selectfont}
\def\supertiny{\fontsize{4pt}{4pt}\selectfont}

\mode<presentation>
{
  \usetheme{Warsaw}
  % \setbeamercovered{transparent}
  \usecolortheme{crane}
}


\usepackage{graphicx, ifthen, listings}

\usepackage[czech]{babel}
% \usefonttheme{professionalfonts}
\usepackage{times}
\usepackage{amsmath}
\usepackage[utf8]{inputenc}
\usepackage{wrapfig}

\usepackage[T1]{fontenc}

\lstset{ basicstyle=\tiny, stringstyle=\ttfamily, showstringspaces=false }

\everymath{\displaystyle}

\setbeamerfont{frametitle}{size=\large}
\setbeamerfont{subsection in toc}{size=\scriptsize}

\makeatletter\newenvironment{blackbox}{%
   \begin{lrbox}{\@tempboxa}\begin{minipage}{0.95\columnwidth}}{\end{minipage}\end{lrbox}%
   \colorbox{black}{\usebox{\@tempboxa}}
}\makeatother

\title[IMF (3)]{Informatika pro moderní fyziky (3)\\rozšíření RubyGems, ladění programů a zpracování textu}

\author[Franti\v{s}ek HAVL\r{U}J, ORF ÚJV Řež]{Franti\v{s}ek HAVL\r{U}J\\{\scriptsize \emph{e-mail: haf@ujv.cz}}}

\date{akademický rok 2019/2020, 9. října 2019}

\institute[ORF ÚJV Řež]
{ÚJV Řež\\oddělení Reaktorové fyziky a podpory palivového cyklu}

\AtBeginSection[]
{
\begin{frame}<beamer>
\frametitle{Obsah}
\tableofcontents[currentsection,hideothersubsections]
\end{frame}
}

\begin{document}

\begin{frame}
  \titlepage
\end{frame}

\begin{frame}
  \tableofcontents
\end{frame}

\section{Co jsme se naučili minule}

\begin{frame}{}
  \begin{itemize}
    \item základy jazyka Ruby na všechny způsoby
    \item vstup a výstup na terminál i do souboru
    \item první kroky ve zpracování dat (čtení ze souboru, zpracování CSV tabulky)
    \item zpracování většího množství dat
  \end{itemize}
\end{frame}
\section{Rozšíření Ruby: RubyGems a Bundler}

\begin{frame}{Knihovny (gemy) jsou základ}
  \begin{itemize}
    \item existují mnohá rozšíření, tzv. knihovny -- v ruby se jim říká \textbf{rubygems}
    \item aktuálně nás zajímá něco na práci s excelovskými soubory
    \item gemy jdou sice instalovat na systémové úrovni, ale z toho je pak zase jenom neštěstí
    \item použijeme radši \textbf{bundler}, správce gemů pro každého: vyřeší za nás závislosti a postará se o snadnou instalaci
  \end{itemize}
\end{frame}

\begin{frame}{Máme bundler?}
  \begin{itemize}
    \item otestujeme rubygems: \texttt{gem -v}
    \item pokud není, zapláčeme, protože jsme asi špatně nainstalovali Ruby
    \item otestujeme bundler: \texttt{bundle -v}
    \item pokud bundler není, doinstalujeme \texttt{gem install bundler}
  \end{itemize}
\end{frame}

\begin{frame}{Jak na to}
  \begin{itemize}
    \item najdu si, která knihovna mě zajímá (třeba na rubygems.org nebo kdekoli jinde): my bychom rádi \textbf{rubyXL} \texttt{https://github.com/weshatheleopard/rubyXL}
    \item vytvořím si prázdný \texttt{Gemfile} -- tam se specifikuje, které gemy chci používat: \texttt{bundle init}
    \item do gemfilu -- je to normální Ruby skript! -- dopíšu\\ \texttt{gem "rubyXL"}
    \item nainstaluju: \texttt{bundle install -{}-path vendor/bundle} (ten parametr stačí poprvé, bundler si to pak pamatuje v konfiguráku \texttt{bundle/.config})
  \end{itemize}
\end{frame}

\begin{frame}{Jak použít?}
  \begin{itemize}
    \item na začátku svého skriptu pak musím nahrát bundler:
    \item \texttt{require "bundler/setup"}
    \item a teď už můžu nahrát jakýkoli gem:
    \item \texttt{require "rubyXL"}
  \end{itemize}
\end{frame}

\subsection{Vytvoříme excelovskou tabulku}

\begin{frame}[fragile]{RTFM, RTFM, RTFM}
  \begin{itemize}
    \item na stránkách \texttt{rubyXL} se nachází spousta příkladů a návodů -- https://github.com/weshatheleopard/rubyXL
    \item kromě toho má i slušnou dokumentaci (GIYF / "rubyxl docs") -- http://www.rubydoc.info/gems/rubyXL/3.3.15
    \item napoprvé navedu do začátku:
  \end{itemize}
  {\scriptsize
  \begin{verbatim}
    workbook = RubyXL::Workbook.new
    worksheet = workbook[0]
    worksheet.add_cell(0, 0, "A1")
    workbook.write("data.xlsx")
  \end{verbatim}
  }
\end{frame}

\begin{frame}{Jednoduché cvičení}
  \begin{itemize}
    \item použijte soubor \texttt{data\_two\_1.csv}
    \item vytvořte excelovský soubor se dvěma listy, na obou bude sloupec 1, sloupec 2 a součet
    \item na jednom součet bude jako číslo (sečte to váš skript)
    \item na druhém bude součet jako excelovský vzorec
  \end{itemize}
\end{frame}


\section{Kde je chyba?}

\begin{frame}{Ladění programů}
  \begin{itemize}
    \item v každém programu je aspoň jedna chyba
    \item není důležité nedělat chyby, ale je nutné je umět najít
    \item když si program/Ruby na něco stěžuje, tak si to přečtěte, jinak se nic nedozvíte
    \item pokud nepoznám, v čem je chyba, jsem bezezbytku ztracen
    \item následují tři úlohy, kde je úkolem najít všechny chyby
    \item z didaktických důvodů postupujte metodou tupého spouštění a postupného opravování
  \end{itemize}
\end{frame}

\section{Zpracování textu}

\subsection{Obecný rozbor}

\begin{frame}{Problém č. 2: mnoho výpočtů, inženýrova smrt}
  \begin{block}{Zadání}
    Při přípravě základního kritického experimentu je pomocí MCNP potřeba najít kritickou polohu regulační tyče R2.

    Jak se tato poloha změní při změně polohy tyče R1?
  \end{block}
\end{frame}

\begin{frame}{Co máme k~dispozici?}
  \begin{block}{MCNP}
    Pokud připravíme vstupní soubor (v~netriviální formě obsahující polohy regulačních tyčí R1 a R2), spočítá nám keff.
  \end{block}
  Potřebovali bychom ale něco na:
  \begin{enumerate}
    \item vytvoření velkého množství vstupních souborů
    \item extrakci keff z~výstupních souborů
    \item popřípadě na vyhodnocení získaných poloh tyčí a keff
  \end{enumerate}
\end{frame}

\begin{frame}{Pracovní postup}
  \begin{enumerate}
    \item načíst keff z~výstupního souboru MCNP
    \pause
    \item vygenerovat potřebné vstupní soubory
    \pause
    \item vyrobit BAT soubor na spuštění výpočtů
    \pause
    \item načíst výsledky ze všech výstupních souborů do jedné tabulky
  \end{enumerate}
\end{frame}

\subsection{Načítání výstupního souboru}

\begin{frame}[fragile]{}
  Nejprve najdeme, kde je ve výstupu z~MCNP žádané keff:
  {\tiny
  \begin{verbatim}
.....
         the k(trk length) cycle values appear normally distributed at the 95 percent confidence level

-----------------------------------------------------------------------------------------------------------------------------------
|                                                                                                                                 |
| the final estimated combined collision/absorption/track-length keff = 1.00353 with an estimated standard deviation of 0.00024   |
|                                                                                                                                 |
| the estimated 68, 95, & 99 percent keff confidence intervals are 1.00329 to 1.00377, 1.00305 to 1.00400, and 1.00289 to 1.00416 |
|                                                                                                                                 |
| the final combined (col/abs/tl) prompt removal lifetime = 1.0017E-04 seconds with an estimated standard deviation of 2.7364E-08 |
.....
  \end{verbatim}
  }
\end{frame}

\begin{frame}{Algoritmus}
  \begin{enumerate}
    \item najít řádek s~keff
    \pause
    \item vytáhnout z~něj keff, takže například:
    \pause
    \item rozdělit podle rovnítka
    \pause
    \item druhou část rozdělit podle mezer
    \pause
    \item vzít první prvek
  \end{enumerate}
\end{frame}


\begin{frame}[fragile]{Realizace (1/5)}
  \scriptsize
  \begin{verbatim}
    keff = nil

    File.foreach("c1_1o") do |line|

    end

    puts keff
  \end{verbatim}
\end{frame}

\begin{frame}[fragile]{Realizace (2/5)}
  \scriptsize
  \begin{verbatim}
    keff = nil

    File.foreach("c1_1o") do |line|
      if line.include?("final estimated combined")

      end
    end

    puts keff
  \end{verbatim}
\end{frame}

\begin{frame}[fragile]{Realizace (3/5)}
  \scriptsize
  \begin{verbatim}
    keff = nil

    File.foreach("c1_1o") do |line|
      if line.include?("final estimated combined")
        a = line.split("=")

      end
    end

    puts keff
  \end{verbatim}
\end{frame}

\begin{frame}[fragile]{Realizace (4/5)}
  \scriptsize
  \begin{verbatim}
    keff = nil

    File.foreach("c1_1o") do |line|
      if line.include?("final estimated combined")
        a = line.split("=")
        b = a[1].split

      end
    end

    puts keff
  \end{verbatim}
\end{frame}

\begin{frame}[fragile]{Realizace (5/5)}
  \scriptsize
  \begin{verbatim}
    keff = nil

    File.foreach("c1_1o") do |line|
      if line.include?("final estimated combined")
        a = line.split("=")
        b = a[1].split
        keff = b[0]
      end
    end

    puts keff
  \end{verbatim}
\end{frame}

\subsection{Zápis všech výsledků do tabulky}

\begin{frame}{Jak na to}
  Máme všechno, co potřebujeme:
  \begin{itemize}
    \item načtení keff z~jednoho výstupního souboru (\texttt{File.foreach}, \texttt{include} a \texttt{split})
    \item procházení adresáře (\texttt{Dir.each})
    \item zápis do souboru (\texttt{File.open} s~parametrem \texttt{w} anebo \texttt{File.write})
  \end{itemize}
  Takže už to stačí jen vhodným způsobem spojit dohromady!
\end{frame}

\begin{frame}[fragile]{Realizace}
  \scriptsize
  \begin{verbatim}
    Dir["*o"].each do |filename|
      keff = nil

      File.foreach(filename) do |line|
        if line.include?("final estimated combined")
          a = line.split("=")
          b = a[1].split
          keff = b[0]
        end
      end

      puts "#{filename} #{keff}"
    end
  \end{verbatim}
\end{frame}

\begin{frame}[fragile]{Výstup}
  Výsledkem je perfektní tabulka:
  {\scriptsize
  \begin{verbatim}
    outputs/c_0_0o 0.94800
    outputs/c_0_10o 0.99800
    outputs/c_0_1o 0.94850
    outputs/c_0_2o 0.95000
    outputs/c_0_3o 0.95250
    outputs/c_0_4o 0.95600
    ...
  \end{verbatim}}
  Hloupé je, že nikde nemáme tu polohu tyčí.
\end{frame}


\begin{frame}[fragile]{Chytrá horákyně}
  ... by jistě vyrobila toto:
  {\scriptsize
  \begin{verbatim}
    0 0 0.94800
    0 10 0.99800
    0 1 0.94850
    0 2 0.95000
    0 3 0.95250
    0 4 0.95600
    ...
  \end{verbatim}
  }
  Nápovědou je funkce \texttt{split} (podle podtržítka) a metoda \texttt{to\_i} (co asi dělá?)
\end{frame}

\begin{frame}[fragile]{Realizace chytré horákyně}
  \scriptsize
  \begin{verbatim}
    Dir["outputs/*o"].each do |filename|
      keff = nil

      File.foreach(filename) do |line|
        if line.include?("final estimated combined")
          a = line.split("=")
          b = a[1].split
          keff = b[0]
        end
      end

      s = filename.split("_")
      puts "#{s[1].to_i} #{s[2].to_i} #{keff}"
    end
  \end{verbatim}
\end{frame}

\section{Automatizace tvorby vstupů}

\begin{frame}[fragile]{Určení poloh tyčí}
  Ve vstupním souboru si najdeme relevantní část:
  \scriptsize
  \begin{verbatim}
    c ---------------------------------
    c polohy tyci (z-plochy)
    c ---------------------------------
    c
    67 pz 47.6000    $ dolni hranice absoberu r1
    68 pz 40.4980    $ dolni hranice hlavice r1
    69 pz 44.8000    $ dolni hranice absoberu r2
    70 pz 37.6980    $ dolni hranice hlavice r2
  \end{verbatim}
\end{frame}

\begin{frame}[fragile]{Výroba šablon}
  Jak dostat polohy tyčí do vstupního souboru? Vyrobíme šablonu, tzn nahradíme
  \begin{verbatim}
    67 pz 47.6000    $ dolni hranice absoberu r1
  \end{verbatim}
  \pause
  nějakou značkou (\emph{placeholder}):
  \begin{verbatim}
    67 pz %r1%    $ dolni hranice absoberu r1
  \end{verbatim}
\end{frame}

\begin{frame}{Chytáky a zádrhele}
  \begin{itemize}
    \item kromě samotné plochy konce absorbéru je nutno správně umístit i z-plochu konce hlavice o 7,102 cm níže
    \item obecně je na místě ohlídat si, že placeholder nebude kolidovat s ničím jiným
  \end{itemize}
  Doporučené nástroje jsou:
  \begin{itemize}
    \item již známá funkce \texttt{sub} pro nahrazení jednoho řetězce jiným
    \item pro pragmatické lenochy funkce \texttt{File.read} načítající celý soubor do řetězce (na což nelze v Pascalu ani pomyslet)
    \item možno ovšem použít i \texttt{File.readlines} (v čem je to lepší?)
  \end{itemize}
\end{frame}

\begin{frame}[fragile]{Realizace}
  \scriptsize
  \begin{verbatim}
    DELTA = 44.8000 - 37.6980

    template = File.read("template")
    (0..10).each do |i1|
      (0..10).each do |i2|
        r1 = i1 * 50
        r2 = i2 * 50
        s = template.sub("%r1%", r1.to_s)
        s = s.sub("%r1_%", (r1 - DELTA).to_s)
        s = s.sub("%r2%", r2.to_s)
        s = s.sub("%r2_%", (r2 - DELTA).to_s)
        File.write("inputs/c_#{i1}_#{i2}", s)
      end
    end
  \end{verbatim}
\end{frame}

\subsection{Co dál?}

\begin{frame}{Navážeme na úspěchy z minulých týdnů}
  \begin{itemize}
    \item vykreslit graf! pro každou z 11 poloh R1 jedna čára (závislost keff na R2)
    \item (= csv soubor, gnuplot, znáte to)
    \item najít automaticky kritickou polohu R2 pro každou z 11 poloh R1
    \item a zase graf... (kritická poloha R2 v závislosti na R1)
  \end{itemize}
\end{frame}

\begin{frame}{A to je vše, přátelé!}
  \begin{center}
    \includegraphics[width=0.8\textwidth]{looney_tunes}
  \end{center}
\end{frame}

\end{document}
